%; whizzy chapter
% -initex iniptex -latex platex -format platex -bibtex jbibtex -fmt fmt
% 以上 whizzytex を使用する場合の設定。

%     Tokyo Debian Meeting resources
%     Copyright (C) 2006 Junichi Uekawa

%     This program is free software; you can redistribute it and/or modify
%     it under the terms of the GNU General Public License as published by
%     the Free Software Foundation; either version 2 of the License, or
%     (at your option) any later version.

%     This program is distributed in the hope that it will be useful,
%     but WITHOUT ANY WARRANTY; without even the implied warranty of
%     MERCHANTABILITY or FITNESS FOR A PARTICULAR PURPOSE.  See the
%     GNU General Public License for more details.

%     You should have received a copy of the GNU General Public License
%     along with this program; if not, write to the Free Software
%     Foundation, Inc., 51 Franklin St, Fifth Floor, Boston, MA  02110-1301 USA

%   Pdf作成手順
% dvipdfmx debianmeetingresume200512.dvi
%  preview (shell-command "xpdf debianmeetingresume200512.pdf&")

%%ここからヘッダ開始。

\documentclass[mingoth,a4paper]{jsarticle}
\usepackage[dvipdfmx]{graphicx}
\usepackage{fancybox}
\usepackage{longtable}
\usepackage{ascmac}	% 囲み (screen,itembox)
\usepackage{fancyvrb}   % 囲み Verbatim のために必要
\usepackage[dvipdfmx]{hyperref}
\usepackage{url}

%http://www.naney.org/diki/dk/hyperref.html
%日本語EUC系環境の時
\AtBeginDvi{\special{pdf:tounicode EUC-UCS2}}
%シフトJIS系環境の時
%\AtBeginDvi{\special{pdf:tounicode 90ms-RKSJ-UCS2}}

%% spacing の設定をする。外枠を減らす。
\setlength\headheight{0mm}
\setlength\topmargin{-20mm}
\setlength\headsep{0mm}
\setlength\topskip{3mm}
\setlength\maxdepth{4pt}
\setlength\columnsep{6mm}
\setlength\textheight{252mm}
\setlength\topmargin{-5mm}
\setlength\textwidth{170mm}
\setlength\oddsidemargin{-5mm}
\setlength\evensidemargin{-5mm}

% commandline環境を定義。画面入出力についてはcommandline環境
% で表記する
\newenvironment{commandline}%
{\VerbatimEnvironment
  \begin{Sbox}\begin{minipage}{15cm}\begin{fontsize}{7.3}{7.3} \begin{BVerbatim}}%
{\end{BVerbatim}\end{fontsize}\end{minipage}\end{Sbox}
  \setlength{\fboxsep}{8pt}\fbox{\TheSbox}}


%%% start of santaku
\makeatletter
\newwrite\tf@jqz
\immediate\openout\tf@jqz\jobname.jqz\relax
\makeatother
\newcounter{santakucounter}
\newcommand{\santaku}[5]{%
\addtocounter{santakucounter}{1}

\addtocontents{jqz}{\arabic{santakucounter}. #5\\}
\nopagebreak 問題\arabic{santakucounter}. 
#1\\
\nopagebreak□ A #2\\
\nopagebreak□ B #3\\
\nopagebreak□ C #4
\pagebreak[1]
\hspace{1cm}
\\

}
%%% end of santaku

\newcommand{\emptyspace}{(\underline{\hspace{1cm}})}

\newcommand{\subsubsubsection}[1]{%
\vspace{1zw}{\bf #1}\\}

% sectionをセンタリングする
\makeatletter
  \renewcommand{\section}{\@startsection{section}{1}{\z@}%
    {\Cvs \@plus.5\Cdp \@minus.2\Cdp}% 前アキ
    {.5\Cvs \@plus.3\Cdp}% 後アキ
    {\normalfont\Large\headfont\raggedright\centering}} % style
\makeatother

% section の代わりの環境
\newcommand{\dancersection}[2]{%
\newpage
東京エリアDebian勉強会 2005
\hrule
\vspace{0.5mm}
\hrule
\hfill{}\includegraphics[width=3cm]{image200502/openlogo-nd.eps}\\
\vspace{-4cm}
\begin{center}
  \section{#1}
\end{center}
\hfill{}#2\hspace{3cm}\space\\
\hrule
\hrule
\vspace{1cm}
}

% for dancerj
\newcommand{\fgref}[1]{図\ref{#1}}
\newcommand{\tbref}[1]{表\ref{#1}}

\begin{document}

\begin{titlepage}

% 毎月変更する部分, 本文の末尾も修正することをわすれずに
\title{
 第11回 東京エリア Debian 勉強会\\事前資料}
\date{2005年12月10日}
\author{Debian勉強会会場係 上川 純一\thanks{Debian Project Official Developer}} 
\maketitle
\thispagestyle{empty}
\end{titlepage}

\newpage
\tableofcontents

\dancersection{Introduction To Debian 勉強会}{上川 純一}

今月のDebian勉強会へようこそ。
これからDebianのあやしい世界に入るという方も、すでにどっぷりとつかってい
るという方も、月に一回Debianについて語りませんか？

目的として下記の二つを考えています。

\begin{itemize}
 \item メールではよみとれない、もしくはよみとってられないような情報を情
       報共有する場をつくる
 \item まとまっていないDebianを利用する際の情報をまとめて、ある程度の塊と
       して出してみる
\end{itemize}

また、東京にはLinuxの勉強会はたくさんありますので、Debianに限定した勉強
会にします。Linuxの基本的な利用方法などが知りたい方は、他でがんばってくださ
い。
Debianの勉強会ということで究極的には参加者全員がDebian Packageを
がりがりと作りながらスーパーハッカーになれるような姿を妄想しています。

Debianをこれからどうするという能動的な展開への土台としての空間を提供し、
情報の共有をしたい、というのが目的です。
次回は違うこと言ってるかもしれませんが、御容赦を。

\subsection{講師紹介}

\begin{itemize}
 \item{上川 純一} 宴会の幹事です。
\end{itemize}

\subsection{事前課題紹介}

今回の事前課題は
「Debianの2005年を振り返って」
というタイトルで200-800文字程度の文章を書いてください。
というものでした。
その課題に対して下記の内容を提出いただきました。

\subsubsection{さわださん}

Debianの2005年を振り返ると言われたら、絶対に被るだろうけどsargeリリースか。w
oodyから3年、ほとんどのパッケージはアップグレードされている、となるとサーバ
が動かなくなるんじゃないかってことでsargeにアップグレードできない人って多い
んじゃないのだろうか？実際、1.xを探している人がいるぐらいだしｗ。解決策とし
ては定期的なリリースとアップグレード時の動作保証、なのだけど、いい意味でも悪
い意味でも「遊びでやっている」Debian Developerの方々にリリーススケジュールや
動作保証を強制するのは難しい気がする。そこらへんはubuntuに頑張ってもらうのが
いいんだろうか。

そういえばDebian GNU/kFreeBSDとかDebian GNU/Solarisとかを頑張っている人もい
たな。UNIX萌えでDebian萌えな私としてはいろんなUNIXでさらにDebianだとうれしい。
常用するかは別として。

日記を振り返ってみると白箱と書いてある。GNU/Linuxでもi386以外の選択肢と言っ
たところか。ボーナスが出たので購入を検討するかな。

\subsubsection{吉田＠板橋さん}

「Sargeがリリースされた」が普通、最大のトピックでしょう。
個人的にはそれにより、woodyからsargeにdist-upgradeしてみました。
いくつか問題も有ったけど最終的には移行完了できました。
移行できたことによるメリット。
セキュリティアップデートに追従できる様になり(特にsamba)。
大抵のパッケージはそこそこ新しい物がapt-getで簡単に入るようになった。
sargeに無いパッケージもsidからのバックポートが簡単になった。
また、「Sarge対応のDebian関係の書籍が出版」された。
個人的にはwoody対応本は内容が古そうなので、ほぼ読まずに
web等を参考資料にしていた。
書籍になって基本的なところが明確になった。
また、パッケージの作り方が明確になり、作りやすくなった。

本題
今年(2005年)のDebian関連のトピックとしては
「東京エリア Debian 勉強会が開催された」。
は外せないでしょう。

で、初参加したところ...なぜか２週間後、
「自分が、勉強会資料の総集編をコミケで売っていた。」
これが個人的に最大のトピック。

\subsubsection{中島　清貴さん}

　今年かなり有名になったのではないだろうか。なにをもって有名というのか分からな
いけれどドラゴン桜ぐらいの知名度があると思う。僕の友人も使いだしたので、いろい
ろ試して教えてもらいたいものだ。ご近所さんで使いだしたのが今年はじめてなので知
っている人が増えて助かった。勉強会のほうは最初の頃かなり参加人数が多かったのに
驚いた。よく定員に入れたものだ。こんなに人間が山手線の外側に集まるものなのか。
あと手が汚れないようにパッケージに入ってるスナック菓子うまい棒をまったく食べな
くなってしまった。Debian Weekly Newsのクイズを全問正解するために食べない。食べ
ないと食べたくなる。そこで食うためにクイズで正解するという戦法だ。しかしまだ全
問正解していないから意味がない。

\subsubsection{岩松 信洋}

\begin{itemize}
 \item \texttt{Debianパッケージメンテナーになった}: 今年は私がDebianに深
く係わりを持った年だった。　今まで使っているだけだったのだが、この勉強会
で知り合った人達の助けを得て、debianパッケージメンテナーになり、開発者側
の立場になったということが今年の私自身の大きな出来事である。　現在NMプロ
セス中で、また先は長いですがDebianに貢献できるデベロッパーになれるよう努
力をしていきたい。
  
 \item \texttt{sarge リリース}: woody から3年。sargeがリリースした。私の
力不足で自分のメンテナンスしているパッケージを取り込めなかったのが悔やま
れます。　時期リリースのetchではこのようなことがないようパッケージメンテ
ナンスをしていきたい。
\end{itemize}

\subsubsection{上川}

2005年、debian sarge が正式にリリースされました。Debian Conference はフィ
ンランドにて開催され、いままでに無い規模のお金が動きながらも、無事に終了
しました。今後継続できるのか、それが一番問題だと思います。Debian の規模
は大きく、期待も大きくなっています。その一方でDebianをささえるインフラは
旧来のままの部分が多いです。この微妙なバランスがどうなるのか、今後目がは
なせないです。日本でのDebian開発者の会の活動についても、低調な感じが否め
ません。今年は特に重要なサーバがダウンしたりしたので、影響が出ました。し
かし、今後、どんどんと日本においてもDebianの存在感が増していくと思われま
すので、なんとかしたいですな。

%%% trivia quiz
\dancersection{Debian Weekly News trivia quiz}{上川 純一}

ところで、Debian Weekly News (DWN)は読んでいますか？
Debian 界隈でおきていることについて書いているDebian Weekly News.
毎回読んでいるといろいろと分かって来ますが、一人で読んでいても、解説が少
ないので、
意味がわからないところもあるかも知れません。みんなでDWNを読んでみましょう。

漫然と読むだけではおもしろくないので、DWNの記事から出題した以下の質問にこたえてみてください。
後で内容は解説します。

\subsection{2005年46号}
\url{http://www.debian.org/News/weekly/2005/46/}
にある11月15日版です。

\santaku
{Debian armebの進捗はどうか}
{やっとgcc/glibc/binutilsが移植された}
{ほとんどのパッケージが移植されている}
{まだ起動もしていない}
{B}

\santaku
{DevJamでJavaの現状について議論があった。その際の認識はどうだったか}
{まだフリーなjavaで全てを実装できていないので、動かないものがある}
{フリーなJavaは充分利用できる状況で、それだけで全てが充足できる。}
{フリーなJavaは全く利用出来ない状態}
{A}

\santaku
{Clam Antivirus について Marc Haberが発表したのは}
{15分毎に更新を確認して、あたらしくなっていたら自動で
volatile.debian.net にアップロードする}
{更新は手動で確認して、メンテナが暇なときにアップデートする。新しいデー
タを常に欲しい人は、頑張って自分でアップデートすること。}
{データ量が多いため、更新はしないので、各自がんばって更新してください。}
{A}

\santaku
{debian-installer etch betaが出ました。Joey Hessがこんなに時間がかかったことに
ついて言明したのは}
{めんどくさかったので放置していたので、こんなに時間がかかりました}
{１０位の項目についてそれぞれで３日づつ遅延要因になるため、一月くらいは
遅れるはめになる}
{ちゃんとハックできる人が参加していないので、コードの品質が下がったため、
こんなに時間がかかりました。}
{B}

\santaku
{SugarCRMはMPL1.1をベースとしたライセンスで配布されている。そのライセン
スはフリーだろうか}
{MPLはMozillaのライセンスなので、その時点でフリーだ}
{ウェブページにフリーソフトだ、と書いてあるので、フリーだ。}
{改変した場合に名前を利用できないことになっているので、名前を変更すれば
よいだろう}
{C}

\santaku
{Debconfの発表資料をDFSGフリーにしようという提案についてAnthony Townsが
した反論は}
{MLでのスレッドなどDFSGフリーでないコンテンツは多数ある。全てがそうである必要はない。}
{ライセンスなんてつけるだけ無駄なので、つけないほうがよいでしょう。}
{あらゆるものはDFSGフリーどころか、全部GPLであるべきなので、GPL以外のラ
イセンスは考えるのもおこがましい。}
{A}

\santaku
{Gabor Gombas さんが、複数の-devパッケージがconflictすることについて苦情
を出した。その対応は}
{-devパッケージがインストールできないのは問題なので、上流のやっている内
容を改変して共存できるようにするのがよい}
{opensslとgnutlsをまぜるほうがライセンス的に適切なので、両方がリンクされ
たパッケージを作る}
{includeファイルのパスなどは開発用のAPIの一部であり、同じパスを利用する
複数の-devパッケージはconflictして当然だ。}
{C}

\santaku
{ping が Linux 専用である点についての議論で、FreeBSDやHurdでも動作させる
ためにパッチを適用することに対してはどういう意見が出たか}
{今後のDebianの一貫性を維持するためにはするべきだ}
{pingなんてBSD上でははやらないのでなくしてもよい}
{あきらかにforkしているため、メンテナンスが大変になる}
{C}

\subsection{2005年47号}
\url{http://www.debian.org/News/weekly/2005/47/}
にある11月22日版です。

\santaku
{Matthias Klose がg++について発表したのは何か}
{g++は今後D言語用のコンパイラによって置き換えられるので、C++なんて古い言
語をつかうのはもうやめろ}
{g++のメモリアロケータが変わるため、またg++ で生成されたライブラリの
ABIが変更になる}
{g++は最適化するために今後はマクロの展開処理を省略する。そのために文法が
若干変更になる}
{B}

\santaku
{Anthony Towns が -private メーリングリストについて提案したのは}
{3年たったら一般公開する}
{存在自体を抹消する}
{即時公開メーリングリストにする}
{A}

\santaku
{Branden RobinsonがDPLについて何ができるかという説明文を発表した。その条
文はいくつあるか}
{3}
{10}
{120}
{B}

\santaku
{Enrico Ziniが発表した新しい検索エンジンでは何をもってパッケージを検索できるか}
{2chの過去ログ情報を用いて検索}
{debtags情報を使って検索}
{popconの利用頻度情報を使って検索}
{B}

\santaku
{Ian Jacksonが提案したのは何か}
{パッケージの自動テストのためのスクリプトインタフェース}
{パッケージを受け入れるときのための基準}
{パッケージの品質をあげるための魔法}
{A}

\santaku
{Christopher Berg が発表した、メンテナ向けのパッケージ一覧ページの新機能
でないのは}
{パッケージがどれくらい人気あるのかということを確認できる}
{パッケージがどれくらいよい品質なのかが確認できる}
{一覧で確認できるパッケージを任意に追加できる}
{B}

\santaku
{PHPライセンスについてSteve Langasek の考えは}
{PHPを使うこと自体がまず問題だ}
{PHP自体については問題ないが、PHP以外にそのライセンスを適用するのには問
題がある}
{PHPライセンスは本当にDFSGフリーなのかどうかはグレーだ}
{B}

\subsection{2005年48号}
\url{http://www.debian.org/News/weekly/2005/48/}
にある11月29日版です。

\santaku
{Freetype に関して何が起きる、とSteve Langasekは宣言したか}
{誰も使っていないので、パッケージを削除する}
{ABIに変更があったので、5のパッケージが移行する必要がある}
{ABIに変更があったので、600のパッケージが移行する必要がある。}
{C}

\santaku
{sbuildの最新版はバージョンが 1.0-1 のパッケージに対してのbinary NMU番号をどうつけてくれるようになったか}
{1.0-1+b1}
{1.0-1.1}
{1.0-1.0.1}
{A}

\santaku
{Frank K\"uster は、パッケージのconffileへの変更の反映を管理者が拒否し、その結
果 postinst が失敗になることについて、問題ないだろう、と質問した。それに
対しての Petter Reinholdtsen の対応は}
{そういうエラーは管理者が拒否するのが問題なので、管理者を日勤教育するべきだ}
{そのような問題は存在しない}
{そういう場合には、設定ファイルを動作に必須なものとローカルで管理者がオー
バライドする部分とに分離することを提案する}
{C}

\santaku
{vserverは何をするものか}
{chrootなどの技術を応用し、複数の仮想サーバコンテキストを作成してくれて、Linux上で複数のサーバを仮
想的に提供できる}
{サインは}
{サーバの統合管理のためのツール}
{A}

\subsection{2005年49号}
\url{http://www.debian.org/News/weekly/2005/49/}
にある12月6日版です。

\santaku
{Manoj SrivastavaがGRの議論期間を宣言した。今回の議論は何についてか}
{-private メーリングリストの一般公開について}
{-devel メーリングリストの秘密化について}
{-mentors メーリングリストの会員制化について} 
{A}

\santaku
{テンポラリディレクトリについての議論があり、ユーザ毎にテンポラリディレ
クトリを持つことがよいのではないかという結論が出た。
ユーザ毎にテンポラリディレクトリを持つ際にその機能を実装してくれるのは}
{/etc/profile でテンポラリディレクトリの作成}
{initスクリプトでのディレクトリの作成}
{pam-tmpdir というPAMモジュール }
{C}

\santaku
{C++のメモリアロケータの移行でまだ移行できていない、ということでさらしあ
げになった日本の開発者は}
{mhattaさんと土屋さん}
{gniibeさんと鵜飼さん}
{えとーさんと岩松さん}
{A}

\santaku
{パッケージがどのバージョン(unstable, stable, testing)用に作成されたのか
を確認する簡単な方法がないか、という質問に対しての Marc Brockschmidt の
回答は何だったか}
{パッケージのバージョン番号を見ればわかる}
{パッケージのchangelogを見ると、どのバージョン用にビルドしたのか、という
ことは確認できる。}
{Debianのパッケージはほとんど全てが一旦はunstableにあったことがあり、
testingとstableに入るため、パッケージがどれ用につくられるというものでは
ない。}
{C}

\dancersection{最近のDebian関連のミーティング報告}{上川 純一}

\subsection{東京エリアDebian勉強会10回目報告}

前回開催した第10回目の勉強会の報告をします。

今回はdpkg-statoverrideとDWNの翻訳についての話を展開しました。
今回の参加人数は7人でした。議論された点を以下に紹介します。
	  
\texttt{dpkg-statoverride} で設定しても、
\texttt{chmod}, \texttt{chown} を\texttt{postinst}で呼んでいる場合は\texttt{dkpg-statoverride}で
指定した値はオーバライドされてしまう。
次に\texttt{dpkg-statoverride}がうごいた場合に直るので、
インストールしてからしばらくすると直る、というような謎のバグの温床になってしまう。


無いユーザ名を指定しての\texttt{statoverride}は
\texttt{「dpkg-statoverride: non-existing user XXX」}というエラーで終了しますね。


\texttt{dpkg-statoverride}コマンドを指定しても、
どのパッケージの\texttt{postinst}で指定されたか、という話がわからない。


現在あまり活用されていないので問題は露見していないが、
\texttt{alternatives}と同じで、\texttt{dpkg}で管理されているわけではないので、
\texttt{dpkg -c} や \texttt{-L} で見えない。
\texttt{alternatives} や \texttt{diversion} と併用された場合の挙動が不明。
おそらくうまくいかないだろうと思われる。この点は確認が必要。

ロギングや、ユーザの作成、削除についてはpostinstで毎回
複雑な処理を追加する必要があり、難しい問題だ。
トレーサビリティーが無い。
ただ、システム管理者側としては、知っていて損はない機構。

\texttt{/usr/lib/dpkg/statoverride.d/}\textit{package名}
というディレクトリがあります、というような仕様のほうがよくないですか？

\texttt{/var/lib/dpkg/info/diversion} のファイル仕様もなかなか微妙だが、
\texttt{/var/lib/dpkg/info/statoverride} のファイル仕様は
ACLやSELINUX対応に拡張できるような形式ではない。
セキュアOS系のタグとかACLとかも管理するために使えそうな仕組みではあるのだが。

DWN翻訳：DWN翻訳については、内容間違ってもユーザのシステムが壊れる、とい
うものではないので、その点は気楽なので、よい。翻訳しにくいときに、参考リ
ンク先の内容をみても全く同じ表現があるときには解決しなくて困る。困った時
には査読の人にお願いするとコメントがくるのでなんとかなることが多い。最近
レビューアとして新しい人が来たのでどうなるか楽しみ。

WML を テキストにして整形するのに \texttt{emacs} の \texttt{auto-fill} 
で毎回処理していたり、\verb![1]!などのURLへのリンクの表現が次の行にうつらないように
禁則処理をする、とかの処理を実は手動で現在している。

対訳表を作成したい。必要。
難しい点としては、よくわからない用語がたくさんある。
また、Debian用語だけをとっても
dpkg、apt,aptitude,dselectなどで用語が統一されていない。
cursesを使っているツールだと、文字数の制限があって、
「インストール」のかわりに「導入」ということばを使っていたりする。
synapticやrpmなどとも用語をそれなりに統一しては行きたいのだが、
訳をしている人達と、話をする機会がないのでなかなかすすまない。
unstable, stable, testingなどを「distribution」というのはちょっと違う気がする。
distributionを配布版と訳するともっと違う気がする。

DWNを翻訳する作業は定期的に行われるのでそれでメンテナンスするとよいものができそうではないか？

DWNに関して、
RSSと、HTML のアンカーが欲しいが、追加するのはWMLのアーキテクチャとして結構難しいかもしれない。
本当にこれでやっているのか、XMLで処理していたりしないのだろうか？

\dancersection{一年間Debian勉強会をやってみて}{上川}
\label{sec:uekawa}
%% 上川の記事はここから 

この記事の目的は、終ってからだと忘れてしまいそうだし、最中だといそがしく
ていっぱいいっぱいなのでどこにも記録されずに忘れ去られてしまいそうな事項
についてメモをしています。

希望としては、ここに書いてある内容をみて、今後のミーティングの運営の手伝
いを人に頼めるようになればよいなと思っています。

\subsection{月例のDebian勉強会のワークフロー}

2005年、Debian勉強会を毎回実施する際に利用したワークフローを紹介します。今後の勉
強会などの参考にできるかと思い、記録します。

参加者規模、10名から20名程度でした。
予算規模は、宴会を含むと一回5万円から10万円程度です。
宴会を含まないのであれば、多くて１万円くらいでした。(\tbref{tab:yosan})

\begin{table}[ht]
 \caption{予算概算}\label{tab:yosan}
 \begin{center}
  \begin{tabular}{|c|c|}
  項目 & 予算 \\
 \hline
  部屋代 & 1500 \\
  コピー代 & $ 300 \times \texttt{人数} $ \\
  宴会代 &  $ 5000 \times \texttt{人数}$ \\
  \end{tabular}
 \end{center}
\end{table}


\subsubsection{1年前}

開催者側のスケジュールの確保。
上川は一年前にだいたいその年のスケジュールを決めています。

\subsubsection{2ヵ月前}

会場の予約確保、開催を決断。

\subsubsection{1ヵ月前}

この時期にすくなくともテーマの設定をします。
講師の確保をしておきます。資料の作成開始をしておかないと間に合わないで
しょう。目処がつきそうだったら、開催のスケジュールを対外的に公表します。
参加者にスケジュールの調整をお願いします。

\subsubsection{1週間前}

宴会の会場選定などを実施します。
大体、資料作成のデッドラインです。
リマインダーの送付をします。

\subsubsection{2日前}

事前課題の文書を事前資料に転記したり、最終的な文書の校正。
この時点で資料の印刷用の最終版が作成。

宴会の人数確定。
宴会予約。

ただ、二日前に選定するとなると場所が限られる場合が多いので、本当はもっと早い時期がよ
いです。一般には、確定が早ければ早いほど予約は安くすみます。
二日前になっても参加できるかどうかわからないという人がいますが、
そういう人の対応は難しいです。
店の柔軟な対応に期待するか、コストをかけるしかないです。
しかし、12月10日の宴会も前日で予約できたのであれば、実は当日に急に開催決
定するとかいうのでさえなければ何とかなる物なのかもしれないです。

\subsubsection{1日前}

資料の印刷をします。
Kinko'sにすべてを依頼する場合は場合にもよりますが、半日くらいは見込む必要があります。
自分で全部するとしても量によりますが、一時間は見込む必要があります。

Kinko's にすべてを依頼する場合、部数が少ないとかなり割高になります。
\footnote{A3の紙にA4を面付けしてもらい、なかとじホッチキス製本にすると
ホッチキスだけで150円/冊になります。コピーが一面14円程度になります。
結果として、一冊450円程度になる。会費を500円しか徴収しないことを考えると、会場費用を
考えると確実に赤字になってしまうので注意。}

\subsubsection{当日}

資料をもっていきます。
司会をします。
適当にもりあがります。

宴会も実施します。
2005年は、講師は無料で宴会、ということで運営しました。
ただ、ときどきそれでは予算が苦しい場合も多々ありました。
なぜか宴会に来ているのに現金をもっていない人とかの扱いには苦慮します。

予算は、ほぼ確実になんらかの理由でのキャンセルが発生するため、
余裕を20％くらい確保できていないと赤字になります。
\footnote{
回避策としては、来ない人から徴収するとかいう案も可能性としてはありますが、
来ない人から徴収するということは暗黙に開催者が次回その来ない人から徴収す
る分について肩代りする、ということを意味するため、オーバヘッドが発生する
ことを忘れてはならない。
}


\subsection{JDMCのような大きなイベントのワークフロー}

Japan Debian Miniconfはまだまだこれから育って行くようなイベントです。
今回蓄積できたノウハウだけで今後もうまく開催できるとは思っていません。
ただ、今回イベントを開催する上で重要でたりなかった点を列記していきます。

\begin{itemize}
 \item 連絡先を明確にする。
 \item 緊急時に判断をできる人を明確にする。
 \item 連絡網を整備する。
 \item ディスカッションができて、そこで決定した事項が合意したとみなせる
       環境をきめてしまう。たとえばIRC。
\end{itemize}

一人ではかぶりきれない責任もあるため、
大きなイベントでは、本気で責任をもって開催したい、と思っている人が複数
いる必要があります。

会議の内容をログに残して全員に周知させる係の人が必要です。
理想としては、実働部隊と分けられればわけたほうがよいです。
JDMCでは、ほとんど矢吹さんだけに情報が集中していたはずで、
メーリングリスト上ではながれていない情報が多数ありました。
もしかすると検討する余裕がなかった項目も多数あったかもしれません。

また、メーリングリストで流れる情報は時系列なので、
現在のステータスを一覧で把握できないです。
タスクトラッキングが重要になります。

また、全員がどういう方法で情報交換をするのかという点について同意が必要で
す。メールで主要な情報交換はなされたのだが、一部の主要メンバーの人達がメー
ルをほぼ全く読んでいなかったという問題がありました。


\subsubsection{2年前}
参加者が稼動できるように日程を確保します。
スポンサーにあたりをつけはじめる。
マネージメント層に交渉します。
それとなく開催できそうな雰囲気がただよっていることを確認します。

\subsubsection{1年前}

一年前か、半年前くらいの時期にスポンサーの予算が大体確定するはずです。
講師に関しての予定、参加者の人数、プログラムの大体のイメージが決まってい
る必要があります。

初の企画でないのなら、前年度のイベントに参加して運営側で何がおきるかを明確にして、
会場のサイジングなどをする必要があります。

スポンサーに関しては、通常スポンサーから資金が提供されるのはイベントが終
了した後です。そのため、事前に当面必要な運転資金をどう確保するのかとい
うのも検討しておく必要があります。

また、赤字になることが見込まれるのであれば、計画を中止するという選択も必
要です。

\subsubsection{6月前}

会場を確保します。
宴会場を確保します。

予約システムを整備し、広報します。
広報は下記を想定しています。

\begin{itemize}
 \item マスメディアへの広報
 \item IRCなどのくちこみ。
       \verb!#debian-devel@opn!など
 \item Blog
 \item DWNへの投稿
 \item メーリングリスト, debian-devel@debian.or.jp,
       debian-users@debian.or.jp, debian-devel@lists.debian.org
 \item Mixiなどのソーシャルネットワーク
 \item Slashdotへ たれこむ
\end{itemize}

また、GPGサイン会などを実施するのなら事前に充分に準備、広報する必要があ
ります。

\subsubsection{1月前}

宴会場の確保、決定が必要です。

参加者の登録が確定しているくらいが本当は好ましいです。
人数が足りないのであればがんばってかきあつめるなどのアクションをとります。

ロジスティックの計画があるので、この時点での人数の把握は重要。

\subsubsection{7日前}

宴会場に連絡して、大体の人数を調整。

\subsubsection{2日前}

宴会場との調整、当日の人数のより確度の高い情報を提供。

\subsubsection{当日}

参加者の出欠確認

参加費用の集金を実施します。

スポンサー企業からの提供物を提供します。スポンサーのグッズとかです。

\subsubsection{事後}

スポンサー企業への報告を作成します。
結果報告書を書き上げます。

参加者の報告をまとめてもらいます。来年のイベントに繋げるために重要です。

次回への検討をはじめます。

\subsubsection{参考文献}

いろいろと他のイベントの報告などもあります。参考になりそうなものを列挙し
ます。

\begin{itemize}
 \item
      Joey の LinuxTagレポート
      \url{http://www.infodrom.org/~joey/Vortraege/2005-06-24/index.html}
 \item
      Joey の LinuxTag感謝状
      \url{http://www.infodrom.org/~joey/log/?200512020951}
 \item
      Debconf5 Final Report
      \url{http://lists.debian.org/debian-devel-announce/2005/12/msg00001.html}
      
\end{itemize}

\subsection{やった内容}

やった内容はけっこういろいろありました。
最初は一般的なうけをねらったものもありましたが、
全体的には技術的な内容を主としています。

\begin{itemize}
 \item 毎月のクイズ
 \item 最初の数回はグループワーク
 \item バックアップリストアについて
 \item ネットワーク監視
 \item reportbugの使い方
 \item debhelper
 \item Social Contract
 \item po-debconf
 \item lintian/linda
 \item dpkg-cross
 \item dsys/update-alternatives
 \item debian-installer
 \item dpatch
 \item toolchain
 \item ITPからアップロードまでの流れ
 \item debconf 2005 参加報告
 \item Debian JP webの改革
 \item debconf の使い方
 \item apt-listbugs
 \item debbugs
 \item dpkg-statoverride
 \item Debian Weekly News 日本語翻訳のフロー
\end{itemize}

来た人数は\tbref{tab:count}にあるような数字です。正確な記録は実は残っていないような気が
しています。議事録をあさればわかるのかもしれません。

\begin{table}[ht]
 \caption{参加人数(概算)}\label{tab:count}
 \begin{center}
  \begin{tabular}{|l|c|}
   & 人数 \\
 \hline
   2005年1月 & 21 \\
   2005年2月 & 10 \\
   2005年3月 (早朝)& 8\\
   2005年4月 & 6\\
   2005年5月 & 8\\
   2005年6月 & 12\\
   2005年7月 & 12\\
   2005年8月 & 7\\
   2005年9月 & 14\\
   2005年10月 & 9\\
   2005年11月 & 8\\
   2005年12月 & 8 \\
  \end{tabular}
 \end{center}
\end{table}


\subsection{おきたトラブル}

勉強会を毎月開催する上で発生したトラブルを紹介します。
\tbref{tab:trouble}です。数字はどれくらいの確率でおきたような気がしてい
るかというのをなんとなく気分的に定量的に書いてみました。

\begin{table}[ht]
 \caption{発生トラブル}\label{tab:trouble}
 \begin{center}
  \begin{tabular}{|l|c|}
   イベント & 発生率 \\
   \hline
   パソコンが盗まれる &  10\%\\
   家が水没する &  10\%\\
   病気で倒れる &  20\%\\
   〆切におくれる &  20\%\\
   なぜか講師のひとと前日まで音信不通 &  10\%\\
   20分くらいまえに連絡してきて、来れないという参加予定者がいる。&  100\%\\
   何も連絡なく来ない人がいる &  100\%\\
   なぜか赤字 &  40\%\\
  \end{tabular}
 \end{center}
\end{table}


\subsection{できた内容}

事前課題により事前にawarenessを向上しました。
いろいろと知らないことを積極的に調べることにより講師がその分野に詳しくな
るという副作用があります。
調査して文章を書いている過程でバグが気に入らないので、バグが直る、という
ことを若干期待しています。

勉強会をクイズではじめてみんなで発言することにより場を和ませることができ
たか？と思っています。クイズは、全員に紙で配布して解いてもらわないと、順
番にあてる形でやると、一部の回答している人だけが集中して、その他の人が当
事者意識をもたないという問題があります(JDMCでの失敗)。
紙を毎回印刷するコストは大きいですが、それなりに効果もあります。

終ってからの blog へのリンク、議事録の掲載についてはあまり反響が無いです。
事前資料のPDFについてはいろいろとblogとかをみているとコメントがあったこ
ともありますが、そちらも反応はあまりないようです。
見られているのかどうか不明です。PDFファイルだからでしょうか？

勉強会の資料を半年分まとめて書籍のような形式にして、
Debian 勉強会資料ということで、コミックマーケットにて販売してもらう、とい
う試みをしています。これは、以前Debian関係の話題が豊富にはいっていた
「Debuan BNU/Linux 不徹底入門」という同人誌があったのですが、それが廃刊
になったため、その代替となれることをめざしているためです。

\subsection{今後やりたいこと}

今後は事前の打合せをもっと密にしたいと考えています。

IRCの debianjp チャンネルで偶然いたメンバーで、なんとなく打合せをする、
ということはできていました。
しかし、最初のころは事実上打合せは上川が電話で呼び出してどっかの飲み屋で
する、という手法をとっていました。
後半は時間の都合で、ほとんど打合せができていなくて、前回の勉強会の後の飲
み会で決定した内容そのままで次の勉強会にのりこむ感じでした。

事後の処理をなんとかしたい、と考えています。
開催した結果をもっと参加していない人にもわかるように効率よくアウトプット
できないだろうか、と思っています。

他の人が参加したいと思えるようなアウトプットが出せないだろうか、と考えて
います。
勉強会自体にDebian関係者が参加したい、と思えるようになることと、
Debianにこれから入る人達が参加したい、と思えるようになることが必要だと思
います。

来年の提案として
システムの構築報告、動作検証、というのはどうだろうか。
「この組合せはできるだろう」、という組合せに関して、
連係はこうやってできる、ということを報告していけば、
多くの人がその動作を確認できるようになり、
問題も解決していけるでしょう。
Debianユーザの勉強会というのはそういう形になるのではないでしょうか。

Debian勉強会以外では、おそらく開発に必要な情報についてまとめて情報収集できる場というのが存在しないため、
開発に必要な情報については継続してやりたいと考えています。
ただ、2回に一回くらいはそういうユーザよりの情報の検証にあててもよいだろ
うと考えています。

また、勉強会でいいっぱなしではなく、勉強会の結果何かが起きる、というようにした
い。
メンテナがバグトラッキングシステムにバグをファイルします、というように宣
言して、毎月その進捗を報告する、という内容にしてみてもよいかな、と思って
います。



\dancersection{Debian勉強会の事前資料の作成はどうやってやったか}{上川}

\subsection{作成ツール}

作成のデータ共有にはaliothのcvsを利用しました。

データの編集は上川はemacs+yatex+whizzytexで実施しました。
\LaTeX 処理系としてplatexを利用しました。
PDFの作成は、dvipdfmxを利用しました。

プリビューはdviファイルはadvi、PDFファイルに関しては、xpdfを利用しました。 

上川の編集環境はDebian sidで、常に開発中の環境だったので、
その時期において動かないツールというのもたまにあり、それなりに大変でした。
原稿の編集中には \texttt{apt-get dist-upgrade}しないように自制していまし
た。

\subsection{\LaTeX ソース}

事前資料は\LaTeX で作成しました。
作業は大きく３種類ありました。

\begin{itemize}
 \item クイズの作成
 \item 参加事前課題の作成
 \item 勉強会のネタの作成
\end{itemize}

\subsubsection{クイズ}
クイズについては、 \LaTeX のマクロでクイズを作成できるようにして、それを利
用して本文を作成しました。

\LaTeX のソースに下記のように記述すると、

\begin{commandline}
 \santaku{問題文}{回答A}{回答B}{回答C}{回答}
\end{commandline}

下記のような出力がでるようになりました。

%\begin{screen}
 \santaku{問題文}{回答A}{回答B}{回答C}{回答}
%\end{screen}

また、その出力を latex-beamer\footnote{\LaTeX でプレゼンテーションを作成す
 るためのスタイル}
 で処理をして、プレゼンテーション形式になるようにしました。
2005年10月以降、勉強会当日は、それを利用して回答を提示するようにしました。

\subsubsection{参加事前課題}

メールにて参加者からplain textできたものを気合いで\LaTeX になおしました。
\LaTeX で使えない文字というのがあるので、それをエスケープすることと、
構造文書については、構造を\LaTeX 用に書き直すという手順が必要です。

例えば、下記のような文章は

\begin{commandline}
□ これについて

こんなことをしてみた

□ あれについて

あんなことをしてみた

□ それについて

いっぱいしてみた

\end{commandline}

itemize環境を利用して下記のような文書になります。

\begin{commandline}
\begin{itemize}
 \item{これについて} こんなことをしてみた
 \item{あれについて} あんなことをしてみた
 \item{それについて} いっぱいしてみた
\end{itemize}
\end{commandline}

\begin{itemize}
 \item{これについて} こんなことをしてみた
 \item{あれについて} あんなことをしてみた
 \item{それについて} いっぱいしてみた
\end{itemize}


\subsubsection{勉強会のネタ}

講師の方に直接 \LaTeX で文書を書いてもらいました。
CVSレポジトリはalioth.debian.orgでホスティングしてもらったので、
そこに共同開発者という形で参加してもらいました。

\LaTeX のスタイルはほぼそのまま jsarticle を採用しています。
ただ、セクションのはじめの部分だけはみかけを派手にしようとして
dancersection というマクロを作って独自に定義しています。
各筆者は dancersection 以下に適当に subsection を作って
文書を作成する、というルールになっています。

\begin{commandline}
\dancersection{一年間Debian勉強会をやってみて}{上川}
\label{sec:uekawa}
%% 上川の記事はここから 
\subsection{セクションの名前 }

文章がだらだらと続く

\subsubsection{セクションの名前 }

.
.
.

\subsection{セクションの名前 }
.
.
.

\end{commandline}

\subsubsection{URLやメールアドレスの処理}

\verb!\url{http://url...}! というように表記しています。
また、メールアドレスも環境を定義するのが面倒なので、そのまま
\verb!\url{メール@アドレス}!という形式にしています。

\subsubsection{特殊文字の処理}

\LaTeX でエスケープが必要な文字については\tbref{tab:tokushu}のように対処しています。
\begin{table}[ht]
 \caption{特殊文字}\label{tab:tokushu}
 \begin{center}
  \begin{tabular}{|c|c|l|}
  文字 & 名称 & 表記 \\
 \hline
 \~{ } & チルダ & \verb!\~{ }! \\
 \underline{ } & アンダーライン & \verb!\underline{ }!\\
 \# &ハッシュ& \verb!\#!\\
 \% &パーセント& \verb!\%!\\
  \end{tabular}
 \end{center}
\end{table}


%% ここまで

\newpage 

\dancersection{次回}{}

未定です。
内容は本日決定予定です。

参加者募集はまた後程。

\newpage

\vspace*{15cm}
\hrule
\vspace{2mm}
\includegraphics[width=2cm]{image200502/openlogo-nd.eps}
\noindent \Large \bf Debian 勉強会資料\\ \\
\noindent \normalfont 2005年12月10日 \hspace{5mm}  初版第1刷発行\\
\noindent \normalfont 東京エリア Debian 勉強会 （編集・印刷・発行）\\
\hrule

\end{document}
